\section{The Folded Napkin: A Claim and a Search}

\subsection{The claim, in its earliest datable words}

The version of the story that circulates has a fixed shape. Here it is as it
was quoted by a reader writing to \emph{Jerusalem Perspective} on 7 October
2006---the earliest dated appearance this study located:

\begin{studyframe}\small
``In order to understand the significance of the folded napkin, you have to
understand a little bit about Hebrew tradition of that day. The folded napkin
had to do with the master and servant, and every Jewish boy knew this
tradition.

When the servant set the dinner table for the master, he made sure that it was
exactly the way the master wanted it \ldots\ If the master was done eating, he
would rise from the table, wipe his fingers, his mouth, and clean his beard,
and would wad up that napkin and toss it onto the table. The servant would then
know to clear the table. For in those days, the wadded napkin meant, `I'm
done.'

But if the master got up from the table, folded his napkin, and laid it beside
his plate, the servant would not dare touch the table because the servant knew
that the folded napkin meant, `I'm not finished yet.' The folded napkin meant,
`I'm coming back!'\,''\endnote{Reader's question printed at the head of
``Christ's Linen Napkin (John 20:7): Is it significant that the napkin \ldots\
was found in the empty tomb folded?'', \emph{JerusalemPerspective.com}, first
published 7 October 2006, read on 25 July 2026 in the Internet Archive capture
of 7 March 2013 at
\texttt{web.archive.org} (capture \texttt{20130307140747} of \texttt{jerusalemperspective.com/3878/}).
The live page was unreachable at that date. The publication date is the site's
own metadata field, and the earliest capture postdates it by six years; this
limit is recorded.}
\end{studyframe}

The claim is admirably specific, and its specificity is what makes it testable.
For it to be true, four things must hold: that table napkins were in domestic
use in first-century Jewish Palestine; that a master's disposal of his napkin
functioned as a signal to a servant; that the signal carried the two contrasting
meanings alleged; and that the Greek word at Jn 20:7 means ``folded.'' If any
one of the four fails, the claim fails.

\subsection{What was searched, and how}

The rabbinic corpus was searched through Sefaria's public text and search
interfaces on 25 July 2026. English phrase searches across the indexed corpus
returned the following, each independently re-run for this study:

\begingroup\small
\noindent\begin{longtable}{>{\raggedright\arraybackslash}p{0.40\linewidth}%
>{\raggedright\arraybackslash}p{0.12\linewidth}%
>{\raggedright\arraybackslash}p{0.40\linewidth}}
\toprule
\textbf{Query} & \textbf{Total hits} & \textbf{Comment}\\
\midrule
\endhead
\texttt{"folded napkin"} & 0 & \\
\texttt{"folds the napkin"} & 0 & \\
\texttt{"folds his napkin"} & 0 & \\
\texttt{"crumpled napkin"} & 0 & \\
\texttt{"napkin" "servant" "signal"} & 0 & \\
\texttt{"sign to his servant"} & 0 & \\
\texttt{"I am coming back"} & 0 & \\
\texttt{"I am not finished"} & 0 & \\
\texttt{soudarion} & 0 & the Greek loanword's own form\\
\texttt{napkin} & 49 & every hit inspected; see below\\
\texttt{sudarium} & 43 & every hit's locus inspected\\
\bottomrule
\end{longtable}
\endgroup

The forty-nine ``napkin'' hits fall into four groups, and not one of them is a
signal: the Mishnah Berakhot 8:3 dispute discussed below and its commentaries;
lexicon entries (Klein, Jastrow) for \emph{mappah}, \emph{mappit},
\emph{mitpahat}; ritual-purity material about cloths and vessels in Kelim,
Chagigah, and Yadayim, and the parallel discussion at Sukkah about a cloth
wrapped round a man; and modern halakhic and responsa literature far outside
the period.

The limits of this search must be stated with the result. Sefaria's index does
not contain every rabbinic work. A phrase search in English cannot find a
custom described in words the searcher did not think of, and it depends on the
translations indexed. What the search establishes is bounded: within the
indexed corpus, in the phrasings tried, there is nothing.

What raises that bounded negative into something stronger is that the places
where such a custom \emph{would} have to be recorded do survive, are indexed,
and were read.

\subsection{The one real napkin text, and what it is about}

There is exactly one ancient rabbinic discussion of a cloth used at table, and
it is the third item in a list of disputes between the House of Shammai and the
House of Hillel ``concerning the meal.'' In the public-domain Hebrew of
Mishnah Berakhot 8:3 the ruling runs, in transliteration: \emph{Beth Shammai
omerim, meqanne'ah yadav ba-mappah u-mannihah al ha-shulhan. U-Veth
Hillel omerim, al ha-keseth.}\endnote{Mishnah Berakhot 8:3, Hebrew, Torat
Emet 357 text, license recorded as public domain, retrieved through the Sefaria
text API on 25 July 2026. Hebrew is transliterated here because this document
carries no Hebrew face; the transliteration is not a translation.}

The Talmud's discussion of it, in A.~Cohen's public-domain 1921 Cambridge
translation, shows what the dispute is for:

\begin{studyframe}\small
``Our Rabbis have taught: Bet Shammai say: One wipes his hands with a napkin
which he lays upon the table; for if thou sayest that [he lays it] upon the
bolster, there is a fear lest the moisture on the napkin may contract
defilement on account of the bolster and this in turn may render the hands
ritually unclean. \ldots\ But Bet Hillel say: [He lays it] upon the bolster;
for if thou sayest that [he lays it] upon the table, there is a fear lest the
moisture on the napkin may contract defilement on account of the table and in
its turn defile the food.''\endnote{Babylonian Talmud, Berakhot 52b; A. Cohen,
\emph{Tractate Berakot}, Cambridge University Press, 1921, retrieved through
the Sefaria text API on 25 July 2026, version license recorded there as public
domain.}
\end{studyframe}

This is the whole of the ancient Jewish law of the table napkin. The question
is where a damp hand-towel may be put so that it does not transmit ritual
impurity. There is no folding and no crumpling. There is no servant. There is
no signal. And the cloth is laid down \emph{during} the meal, after the
hand-washing, not at its end.

\subsection{What actually signalled that a meal was over}

The rabbis did have a marker for the end of a meal, and it was legally
operative, which is why it is discussed. It was the removal of the table.
Cohen's translation of Berakhot 42a:

\begin{studyframe}\small
``Rab Pappa went on a visit to the house of Rab Huna b.\ Rab Nathan. After they
had finished their meal, they set before them several things to eat. Rab Pappa
took of them and ate. They said to him, `Does not the master hold that if one
has finished [his meal], he must not eat anything further?' He replied, `This
has only been said when the table has been cleared.'\,''\endnote{Babylonian
Talmud, Berakhot 42a; Cohen, \emph{Tractate Berakot} (1921), via the Sefaria
text API, 25 July 2026.}
\end{studyframe}

This is the decisive negative. Here is the exact legal question the alleged
custom would answer---how does anyone know the meal is over?---being answered,
in a household with servants, by a physical act; and the act is the removal of
the table, not a gesture with a cloth. A custom of the kind alleged, if it had
existed, would have had to be reckoned with on this page. It is not there.

\subsection{A real \emph{sudar} signal, and what it signalled}

The search did turn up one genuine, ancient, Jewish use of a \emph{sudar}---the
same word, in its Aramaic-Hebrew form, that Greek borrowed as
\emph{soudarion}---as a deliberate signal. It is worth reporting, both because
it is a positive finding and because it shows that the corpus is perfectly
capable of recording such a thing when one exists.

Mishnah Sanhedrin 6:1 describes the procession to a stoning:

\begin{studyframe}\small
``One would stand before the entrance to the courthouse, his scarf in his hand;
and another would ride on a horse before them at a distance, in order to see
him. [If] one says, `I have an argument for acquittal,' then he waves the
scarf, and the horse runs to stop him.''\endnote{Mishnah Sanhedrin 6:1,
Sefaria Community Translation, license recorded as CC0, retrieved 25 July 2026.
The Hebrew, in the public-domain Torat Emet text, reads \emph{ehad omed al
petah beth din veha-sudarin beyado \ldots\ hallaz meniph ba-sudarin
veha-sus rats u-ma'amido}.}
\end{studyframe}

So the corpus knows a \emph{sudarin} used as a signal. The signal is judicial,
public, and made by waving; its meaning is ``stop the execution''; and it has
nothing whatever to do with a table, a master, a servant, folding, or a
promise to return. A search capable of finding this would have found the
alleged dining custom if the corpus contained it.

\subsection{The Greco-Roman side}

The claim is about Jewish custom, but the Roman world is where table napkins as
such are best documented, and it should be checked for completeness.

The Roman dinner napkin, \emph{mappa}, was a personal possession that guests
brought to the table and took home again---in Petronius the guests fill their
napkins with fruit, and one of them has apples tied up in his to take back to
his slave. Smith's \emph{Dictionary of Greek and Roman Antiquities} (1875)
describes the \emph{mappa} as ``a smaller kind of napkin, or a handkerchief,
which the guests carried with them to the table,'' and its article on the
subject---which covers materials, sizes, and even the wearing of napkins as
head-dress---says nothing about any signalling function at
table.\endnote{William Smith, ed., \emph{A Dictionary of Greek and Roman
Antiquities} (1875), s.v.\ ``Mantele,'' read on 25 July 2026 at the
LacusCurtius transcription,
\texttt{penelope.uchicago.edu}. Read in an electronic transcription of the
printing, not in a printed volume; that ceiling is recorded.} The same article
notes in passing that at a Greek or Roman meal it was the boy or slave pouring
the water who ``also held the napkin or towel for wiping the hands dry''---so
that the cloth the claim imagines a master manipulating was, in the setting
where such cloths are best documented, in the servant's hands already. An object
each diner brings and takes away is in any case the least available thing at a
Roman table to serve as the host's signal.

There is one famous ancient napkin signal, and its sense is the reverse of the
claim's. Suetonius reports that Nero gave the public the chance to watch him
drive in the Circus Maximus, ``one of his freedmen dropping the napkin from the
place usually occupied by the magistrates.''\endnote{Suetonius, \emph{Nero}
22.2; J.~C. Rolfe's public-domain Loeb translation, read on 25 July 2026 at the
LacusCurtius transcription,
\texttt{penelope.uchicago.edu}.}
The dropped \emph{mappa} started the races. It meant ``go,'' not ``I am
returning.''

\subsection{Lightfoot's silence}

One more negative deserves separate statement, because of who produced it.

John Lightfoot's \emph{Horae Hebraicae et Talmudicae} is the seventeenth
century's great attempt to explain the Gospels out of rabbinic literature. It
proceeds verse by verse and cites Talmudic material relentlessly; on John 20
alone Lightfoot pauses to quote Bava Bathra on the depth of burial niches at
verse 5. If a rabbinic custom of the table napkin bore on John 20:7, Lightfoot
is the man in the history of scholarship most likely to have said so.

All four volumes of Robert Gandell's Oxford edition of 1859 were retrieved and
searched for this study. The word ``napkin'' occurs twice in the whole work,
both times in the note on John 11:44 and both times about the Lazarus miracle
rather than any custom: ``The dead man came forth, though bound hand and foot
with his graveclothes, and blinded with the napkin.'' And on John 20 Lightfoot
comments on verses 1, 5, 12, 17, 23, 24, 25, 26, and 29. He has no note on
verse 7 at all.\endnote{John Lightfoot, \emph{Horae Hebraicae et Talmudicae:
Hebrew and Talmudical Exercitations upon the Gospels}, ed.\ Robert Gandell
(Oxford: University Press, 1859), vol.~3, the John exercitations, chapter XX;
the note on Jn 11:44 is at vol.~3. Full text of all four volumes retrieved from
the Internet Archive and searched on 25 July 2026. This is an OCR search of a
scanned text; an OCR miss cannot be excluded, but the chapter's own sequence of
verse lemmata was read directly and contains no verse 7.}

\subsection{The Greek, once more, and the English splice}

The fourth condition the claim requires is that \emph{entetyligmenon} means
``folded.'' Section~3 established that it does not: Thayer glosses it at this
verse ``to roll up, wrap together,'' and its only other New Testament uses
describe the wrapping of a corpse in a shroud.

There is a further point about English that explains a great deal. The word
``napkin'' at Jn 20:7 comes from the Authorized Version and the Douay, both of
which use an archaic English word for a wiping-cloth. But the Authorized
Version does not say ``folded''---it says ``wrapped together in a place by
itself.'' The word ``folded'' comes from twentieth-century versions. The claim
therefore requires the noun from one translation and the verb from another, and
holds together in no single English Bible.

David Bivin, answering the reader in 2006, reached the same conclusion from the
other end: he noted that the case requires all of ``if first-century Jewish
residents of the land of Israel used table napkins, and if there were such a
custom as described, and if the handkerchief mentioned in John 20:7 were a
table napkin, and if the Greek word \emph{entetyligmenon} meant `having been
folded' rather than `having been wrapped up','' and guessed that ``the detailed
description of this supposed custom is an invention triggered in someone's
fertile mind by the archaic KJV translation, `napkin'.'' He added that no
version of the claim ``cites a biblical or rabbinic source.'' Bivin also
records a smaller negative worth keeping: ``Apparently, there is no early
rabbinic source that discusses how the hands were dried after washing
them.''\endnote{David Bivin's response at \emph{JerusalemPerspective.com}, as
cited above. Quoted in short clauses with attribution; the page is under
present copyright and is not reproduced beyond the argument's need.}

\subsection{Where the claim came from}

The following account is offered as this study's own reconstruction, bounded by
what could be reached.

The earliest datable appearance is the October 2006 \emph{Jerusalem
Perspective} page, and what appears there is not the origin but a
reader's report of something already circulating: the claim arrives already in
its finished form, in a reader's question, and is refuted on the spot by a
specialist in the Jewish background of the Gospels. Its medium was the
forwarded email; the earliest independently archived instance located is a
mailing-list forward of 21 May 2007 carrying the now-standard wording.

A correction to the received account should be recorded. The date usually given
for the claim's origin is 2007, and it traces to a single sentence in
TruthOrFiction's assessment of January 2008: ``The only references to this story
that we found are from Internet postings and emails that seem to have originated
in 2007.'' That sentence was copied, nearly verbatim, into GotQuestions'
treatment---``The only references to this story seem to be from internet
postings and emails that appear to have originated in 2007''---and from there
into a good deal of Catholic and Evangelical apologetic writing. It was a guess,
honestly labelled as one by its author, and the \emph{Jerusalem Perspective}
page pushes the claim back at least to October 2006. Whether it circulated in
print earlier could not be established: Google Books, HathiTrust, and the
Usenet archives were unreachable during this research, and that limit is
recorded here rather than papered over.

What can be said with confidence is the thing that matters. Across every
instance located---the 2006 and 2007 texts, the assessments at TruthOrFiction
and GotQuestions, and the devotional and newspaper retellings that continue to
appear---no version of the claim has ever cited a Jewish source. Its entire
evidentiary apparatus consists of the assertion that ``every Jewish boy knew
this tradition.''

\subsection{The finding}

\begin{studybox}{Project synthesis: the master-and-servant napkin custom}
\textbf{Judgment.} The custom is unattested. No Mishnaic, Talmudic, Toseftan,
midrashic, or minor-tractate text known to this study describes a napkin
gesture as a signal from a master to a servant, or attaches the meanings ``I am
finished'' and ``I am coming back'' to any disposal of a cloth. The claim is a
modern devotional legend, in circulation by 2006, resting on an archaic English
noun from one Bible version and a mistranslated verb from another.

\textbf{Reasons.} (1) Targeted searches of the indexed rabbinic corpus for the
claim's own vocabulary return nothing. (2) The one ancient rabbinic discussion
of a napkin at a meal, Mishnah Berakhot 8:3 with Berakhot 52b, is a
ritual-purity dispute about where a damp hand-towel may be laid; it has no
servant, no signal, and no folding. (3) The rabbis' actual end-of-meal marker,
discussed at Berakhot 42a in exactly the household context the claim imagines,
is the removal of the table. (4) The corpus does record a \emph{sudar} used as
a deliberate signal---at Mishnah Sanhedrin 6:1, to halt an execution---which
shows that the search would have found a dining signal had one existed. (5)
Lightfoot, hunting rabbinic parallels through John verse by verse, has no note
on Jn 20:7 and no napkin custom anywhere in four volumes. (6) The Greek
participle means wrapped or rolled, not folded, and is the verb used of
swathing a corpse. (7) No proponent of the claim, in any instance located from
2006 to the present, has ever cited a source.

\textbf{The strongest counterargument, at full strength.} Absence of evidence in
a literature that is fragmentary, selectively preserved, and concerned with law
rather than with manners is genuinely weak evidence of absence. The rabbinic
corpus records what was disputed, and an uncontroversial domestic gesture would
generate no dispute and therefore no text. Ordinary table manners of the sort
everyone knows are exactly what such literature does \emph{not} record: the
Mishnah preserves the Shammai--Hillel argument about where the napkin goes
precisely because it was contested, and passes over in silence a hundred things
about the same meal that were not. Furthermore the searches performed here were
in English, over an incomplete index, in phrasings a modern researcher chose.

\textbf{Why the counterargument does not carry.} It would carry against a
merely quiet negative. It does not carry against this one, for three reasons.
First, the claim is not about a quiet gesture: it asserts a codified signalling
convention with two contrasting meanings, operative between masters and
servants, that ``every Jewish boy'' knew---something on that scale leaves
traces. Second, the corpus does address the exact governing question, at
Berakhot 42a, and answers it with a different mechanism; a rival signal cannot
hide behind silence on a page where the question is being decided. Third, the
claim's own defenders have never produced a text, and the absence has been
publicly noted by specialists since 2006 without any source emerging.

\textbf{What would change this judgment.} A single ancient Jewish text---
tannaitic, amoraic, midrashic, or from the minor tractates---in which the
disposal, folding, or crumpling of a cloth at table communicates anything to a
servant about whether the meal is over. Failing that, a Second Temple or
Greco-Roman source attesting a napkin signal with the alleged meanings would
substantially weaken the negative even if it were not Jewish. Nothing less
would do; a modern author repeating the claim, however respectable, adds
nothing, since the whole difficulty is that the modern authors have no source
either.
\end{studybox}

\subsection{A note on tone}

It should be said plainly that the people who repeat this story are not being
dishonest, and that what they are trying to say with it is true. Christ is
coming back. The Church has always said so, and says so in the Creed. The
objection to the folded-napkin story is not that its conclusion is false but
that its premise is invented, and that a true doctrine defended by a
manufactured custom is worse defended than not defended at all---because the
person who eventually discovers that no Jewish boy ever knew any such tradition
is left wondering what else was made up.

There is also a loss of a different kind. The verse has a real strangeness in
it, attested in every manuscript and every version, which the invented custom
covers over: a burial cloth still in the shape a head gave it, lying by itself
in a place where a body is not. That is stranger than the legend, and it is
actually there.
